% College Lecture 2 (extended Euclid: process, algorithm, 161/28, 17/0, 0/45, 84/320) + Lecture 4 (inverse with the t columns,
% 11/26, 23/100, 12/26; linear Diophantine equations), Forouzan 2.1--2.2. All tables generated by script (scratchpad/r4/verify.py).
\section{Extended Euclidean Algorithm \hothot}\label{sec:exteuclid}
\asked{Mid-sem 2024--25 Q7 (5 marks): pseudocode + $23^{-1}$ in $\mathbb{Z}_{100}$ with ``value in iteration $k$'' sub-questions. 2025--26 Q5 (3 marks): solve $19x + 13y = 20$. The same table, with polynomials, gives the GF$(2^8)$ inverse asked in both years (\S\ref{sec:algebra}), and $d$ in RSA (\S\ref{sec:rsa}).}

Euclid's algorithm decides whether $\gcd(a, n) = 1$, and hence whether $a^{-1} \bmod n$ exists, but it does not produce the inverse. The extended Euclidean algorithm does. \imp{Given two integers $a$ and $b$, the extended Euclidean algorithm computes $\gcd(a, b)$ and, at the same time, integers $s$ and $t$ such that $s \times a + t \times b = \gcd(a, b)$.} That such $s$ and $t$ always exist is \textbf{Bézout's identity}; the algorithm proves it by constructing them. One table then answers four kinds of exam question: the inverse $b^{-1} \bmod n$, the solutions of a linear Diophantine equation $ax + by = c$, the solutions of a linear congruence (\S\ref{sec:modular}), and the RSA private key $d = e^{-1} \bmod \phi(n)$. With polynomials instead of integers the same table gives inverses in GF$(2^8)$.

\subsection{The idea: carry every remainder as a combination of \texorpdfstring{$a$ and $b$}{a and b}}
By hand, $s$ and $t$ can be found by back-substitution. For $a = 19$ and $b = 13$, Euclid's divisions are $19 = 1 \times 13 + 6$ and $13 = 2 \times 6 + 1$. Solving them from the last line upwards gives
\[
1 = 13 - 2 \times 6 = 13 - 2 \times (19 - 1 \times 13) = 3 \times 13 - 2 \times 19,
\]
so $s = -2$ and $t = 3$. This works, but it requires storing every quotient, running the divisions a second time in reverse, and keeping track of signs, which is where most errors happen. The extended algorithm does the same bookkeeping forwards, in a single pass.

The key observation is that every remainder in Euclid's sequence is an integer combination of $a$ and $b$. The first two numbers of the sequence are $a$ and $b$ themselves, and trivially
\[
r_1 = a = 1 \cdot a + 0 \cdot b, \qquad r_2 = b = 0 \cdot a + 1 \cdot b,
\]
which gives the starting values $s_1 = 1$, $s_2 = 0$, $t_1 = 0$, $t_2 = 1$. Suppose that in some row the \textbf{invariant} $r_1 = s_1 a + t_1 b$ and $r_2 = s_2 a + t_2 b$ holds. The next remainder is $r = r_1 - q r_2$, and substituting the two combinations gives
\[
r = (s_1 a + t_1 b) - q\,(s_2 a + t_2 b) = (s_1 - q s_2)\,a + (t_1 - q t_2)\,b .
\]
So if $s$ and $t$ are updated by exactly the same rule as $r$, namely $s = s_1 - q s_2$ and $t = t_1 - q t_2$, the new remainder satisfies $r = s a + t b$. After the shift $(r_1, r_2) \leftarrow (r_2, r)$, $(s_1, s_2) \leftarrow (s_2, s)$, $(t_1, t_2) \leftarrow (t_2, t)$ the invariant holds again, so by induction it holds in every row. The $r$ columns are Euclid's algorithm unchanged, so the loop ends with $r_1 = \gcd(a, b)$ and $r_2 = 0$, and the invariant for $r_1$ then reads $\gcd(a, b) = s_1 a + t_1 b$. This proves Bézout's identity for all $a$ and $b$.

\subsection{The algorithm and the table (Forouzan)}
The slides give the algorithm as three identical update blocks, one for each pair of variables:
\begin{center}
\begin{minipage}{0.62\textwidth}
\small\ttfamily
r$_1$ $\leftarrow$ a;\quad r$_2$ $\leftarrow$ b;\\
s$_1$ $\leftarrow$ 1;\quad s$_2$ $\leftarrow$ 0;\hfill\textrm{(initialisation)}\\
t$_1$ $\leftarrow$ 0;\quad t$_2$ $\leftarrow$ 1;\\
while (r$_2$ > 0)\\
\{\\
\hspace*{1.5em}q $\leftarrow$ r$_1$ / r$_2$;\\
\hspace*{1.5em}r $\leftarrow$ r$_1$ $-$ q $\times$ r$_2$;\quad r$_1$ $\leftarrow$ r$_2$;\quad r$_2$ $\leftarrow$ r;\hfill\textrm{(updating r's)}\\
\hspace*{1.5em}s $\leftarrow$ s$_1$ $-$ q $\times$ s$_2$;\quad s$_1$ $\leftarrow$ s$_2$;\quad s$_2$ $\leftarrow$ s;\hfill\textrm{(updating s's)}\\
\hspace*{1.5em}t $\leftarrow$ t$_1$ $-$ q $\times$ t$_2$;\quad t$_1$ $\leftarrow$ t$_2$;\quad t$_2$ $\leftarrow$ t;\hfill\textrm{(updating t's)}\\
\}\\
gcd(a, b) $\leftarrow$ r$_1$;\quad s $\leftarrow$ s$_1$;\quad t $\leftarrow$ t$_1$
\end{minipage}
\end{center}
On paper each pass of the loop is one row of a table with the columns $q \mid r_1\ r_2\ r \mid s_1\ s_2\ s \mid t_1\ t_2\ t$. In every row the three blocks do the same thing: the new value is (first) $-\ q \times$ (second), and in the next row the second value moves into the first column and the new value into the second. A final row, with no quotient, shows the values after the loop, and the answer is read from its \emph{first} columns. The slide example, $a = 161$, $b = 28$:
\begin{center}\small
\begin{tabular}{@{}r|rrr|rrr|rrr@{}}
\toprule
$q$ & $r_1$ & $r_2$ & $r$ & $s_1$ & $s_2$ & $s$ & $t_1$ & $t_2$ & $t$\\
\midrule
5 & 161 & 28 & 21 & 1 & 0 & 1 & 0 & 1 & $-5$\\
1 & 28 & 21 & 7 & 0 & 1 & $-1$ & 1 & $-5$ & 6\\
3 & 21 & 7 & 0 & 1 & $-1$ & 4 & $-5$ & 6 & $-23$\\
 & \textbf{7} & 0 & & $\mathbf{-1}$ & 4 & & \textbf{6} & $-23$ & \\
\bottomrule
\end{tabular}
\end{center}
$\gcd(161, 28) = 7$, $s = -1$, $t = 6$: \;$(-1)(161) + 6(28) = -161 + 168 = 7$ \checkmark. The invariant can be checked in every row: $21 = 1 \times 161 - 5 \times 28$, $7 = -1 \times 161 + 6 \times 28$ and $0 = 4 \times 161 - 23 \times 28$. The last of these gives a free self-check: in the final row $s_2 a + t_2 b = 0$, and $(s_2, t_2)$ always equals $\pm(b/d, -a/d)$ with $d = \gcd(a, b)$; here $(4, -23) = (28/7, -161/7)$. In an inverse table (next subsection) this means the final $t_2$ is $\pm n$, so any other value there reveals an arithmetic slip.

The other slide examples test the edge cases. For $a = 17$, $b = 0$ the loop never runs, because $r_2 = 0$ from the start, and the initial values give $\gcd(17, 0) = 17$, $s = 1$, $t = 0$. For $a = 0$, $b = 45$ one iteration with $q = 0$ swaps the pair, giving $\gcd(0, 45) = 45$, $s = 0$, $t = 1$. The slide exercise $a = 84$, $b = 320$ also starts with the smaller number, and the first row ($q = 0$) simply swaps it:
\begin{center}\small
\begin{tabular}{@{}r|rrr|rrr|rrr@{}}
\toprule
$q$ & $r_1$ & $r_2$ & $r$ & $s_1$ & $s_2$ & $s$ & $t_1$ & $t_2$ & $t$\\
\midrule
0 & 84 & 320 & 84 & 1 & 0 & 1 & 0 & 1 & 0\\
3 & 320 & 84 & 68 & 0 & 1 & $-3$ & 1 & 0 & 1\\
1 & 84 & 68 & 16 & 1 & $-3$ & 4 & 0 & 1 & $-1$\\
4 & 68 & 16 & 4 & $-3$ & 4 & $-19$ & 1 & $-1$ & 5\\
4 & 16 & 4 & 0 & 4 & $-19$ & 80 & $-1$ & 5 & $-21$\\
 & \textbf{4} & 0 & & $\mathbf{-19}$ & 80 & & \textbf{5} & $-21$ & \\
\bottomrule
\end{tabular}
\end{center}
$\gcd(84, 320) = 4$, $s = -19$, $t = 5$: \;$-19 \times 84 + 5 \times 320 = -1596 + 1600 = 4$ \checkmark, as on the slide. The final $(s_2, t_2) = (80, -21) = (320/4, -84/4)$ confirms it.

\subsection{Multiplicative inverse with the table \hothot}
\likely{``Find $b^{-1}$ in $\mathbb{Z}_n$ using the extended Euclidean algorithm; give the number of iterations / $r$ or $t$ in iteration $k$'' (5 marks), and $d = e^{-1} \bmod \phi(n)$ in every RSA question.}
To find $b^{-1} \bmod n$, run the algorithm with $a = n$, that is, $r_1 = n$ and $r_2 = b$. At the end $s \times n + t \times b = \gcd(n, b)$. If the gcd is 1, reduce this equation modulo $n$; the slides do it in three lines:
\[
(s \times n + b \times t) \bmod n = 1 \;\Rightarrow\; [(s \times n) \bmod n] + [(b \times t) \bmod n] = 1 \;\Rightarrow\; (b \times t) \bmod n = 1 .
\]
So \imp{the multiplicative inverse of $b$ is the final $t_1$, mapped into $\mathbb{Z}_n$}. The $s$ columns are never used and are dropped, leaving $q \mid r_1\ r_2\ r \mid t_1\ t_2\ t$. If the final $r_1$ is not 1, then $\gcd(n, b) > 1$ and no inverse exists (\S\ref{sec:modular}); the algorithm detects this by itself. A negative $t_1$ is normal and is mapped into $\mathbb{Z}_n$ by adding $n$. The slide's version of the algorithm, which is the pseudocode asked in 2024--25 Q7:
\begin{center}
\begin{minipage}{0.56\textwidth}
\small\ttfamily
r$_1$ $\leftarrow$ n;\quad r$_2$ $\leftarrow$ b;\quad t$_1$ $\leftarrow$ 0;\quad t$_2$ $\leftarrow$ 1;\\
while (r$_2$ > 0)\\
\{\\
\hspace*{1.5em}q $\leftarrow$ r$_1$ / r$_2$;\\
\hspace*{1.5em}r $\leftarrow$ r$_1$ $-$ q $\times$ r$_2$;\quad r$_1$ $\leftarrow$ r$_2$;\quad r$_2$ $\leftarrow$ r;\\
\hspace*{1.5em}t $\leftarrow$ t$_1$ $-$ q $\times$ t$_2$;\quad t$_1$ $\leftarrow$ t$_2$;\quad t$_2$ $\leftarrow$ t;\\
\}\\
if (r$_1$ = 1) then b$^{-1}$ $\leftarrow$ t$_1$ mod n\\
else ``no inverse''
\end{minipage}
\end{center}
The three slide examples:
\begin{center}\small
\begin{tabular}{@{}r|rrr|rrr@{}}
\multicolumn{7}{c}{$11^{-1} \bmod 26$}\\
\toprule
$q$ & $r_1$ & $r_2$ & $r$ & $t_1$ & $t_2$ & $t$\\
\midrule
2 & 26 & 11 & 4 & 0 & 1 & $-2$\\
2 & 11 & 4 & 3 & 1 & $-2$ & 5\\
1 & 4 & 3 & 1 & $-2$ & 5 & $-7$\\
3 & 3 & 1 & 0 & 5 & $-7$ & 26\\
 & \textbf{1} & 0 & & $\mathbf{-7}$ & 26 & \\
\bottomrule
\end{tabular}
\quad
\begin{tabular}{@{}r|rrr|rrr@{}}
\multicolumn{7}{c}{$23^{-1} \bmod 100$}\\
\toprule
$q$ & $r_1$ & $r_2$ & $r$ & $t_1$ & $t_2$ & $t$\\
\midrule
4 & 100 & 23 & 8 & 0 & 1 & $-4$\\
2 & 23 & 8 & 7 & 1 & $-4$ & 9\\
1 & 8 & 7 & 1 & $-4$ & 9 & $-13$\\
7 & 7 & 1 & 0 & 9 & $-13$ & 100\\
 & \textbf{1} & 0 & & $\mathbf{-13}$ & 100 & \\
\bottomrule
\end{tabular}
\quad
\begin{tabular}{@{}r|rrr|rrr@{}}
\multicolumn{7}{c}{$12^{-1} \bmod 26$}\\
\toprule
$q$ & $r_1$ & $r_2$ & $r$ & $t_1$ & $t_2$ & $t$\\
\midrule
2 & 26 & 12 & 2 & 0 & 1 & $-2$\\
6 & 12 & 2 & 0 & 1 & $-2$ & 13\\
 & \textbf{2} & 0 & & $-2$ & 13 & \\
\bottomrule
\end{tabular}
\end{center}
$11^{-1} \equiv -7 \equiv \mathbf{19} \pmod{26}$ (check: $11 \times 19 = 209 = 8 \times 26 + 1$); \;$23^{-1} \equiv -13 \equiv \mathbf{87} \pmod{100}$ (check: $23 \times 87 = 2001$); \;$12^{-1} \bmod 26$: the final $r_1 = 2$, so $\gcd(26, 12) = 2$ and \textbf{no inverse exists}. In the third table $t_1 = -2$ still satisfies $12 \times (-2) \equiv 2 \pmod{26}$, the gcd rather than 1, which is exactly what the invariant predicts.

\begin{worked}[: 2024--25 Q7 (5 marks), inverse with ``iteration $k$'' sub-questions]
\textbf{Q.} Write the pseudocode for finding the multiplicative inverse of $b$ in $\mathbb{Z}_n$, and use it to find $23^{-1}$ in $\mathbb{Z}_{100}$. State (a) the number of iterations, (b) the remainder in iteration 2, (c) the value of $t$ in iteration 4, (d) the iteration whose $t$ becomes the inverse.

\textbf{A.} Pseudocode: as above ($r_1 = n$, $r_2 = b$, $t_1 = 0$, $t_2 = 1$; loop while $r_2 > 0$; inverse $= t_1 \bmod n$ if $r_1 = 1$). An \emph{iteration} is one pass of the while loop, i.e.\ one row that has a quotient; the row shows the values at the start of the pass and the $r$, $t$ computed in it.
\begin{center}\small
\begin{tabular}{@{}c|c|rrr|rrr@{}}
\toprule
Iteration & $q$ & $r_1$ & $r_2$ & $r$ & $t_1$ & $t_2$ & $t$\\
\midrule
1 & 4 & 100 & 23 & 8 & 0 & 1 & $-4$\\
2 & 2 & 23 & 8 & \textbf{7} & 1 & $-4$ & 9\\
3 & 1 & 8 & 7 & 1 & $-4$ & 9 & $\mathbf{-13}$\\
4 & 7 & 7 & 1 & 0 & 9 & $-13$ & \textbf{100}\\
(end) & & \textbf{1} & 0 & & $\mathbf{-13}$ & 100 & \\
\bottomrule
\end{tabular}
\end{center}
Since $r_1 = 1$, $\gcd(100, 23) = 1$ and the inverse exists. (a) \textbf{4 iterations}. (b) The remainder in iteration 2 is $\mathbf{7}$ ($23 = 2 \times 8 + 7$). (c) $t$ in iteration 4 is $\mathbf{100}$, which is $\pm n$ as the self-check predicts. (d) Iteration 3 computes $t = -13$, which becomes the final $t_1$; mapped into $\mathbb{Z}_{100}$: $23^{-1} = -13 + 100 = \mathbf{87}$.
\textbf{Check:} $23 \times 87 = 2001 = 20 \times 100 + 1$ \checkmark.
\end{worked}

\begin{worked}[: inverse in $\mathbb{Z}_{180}$ (Forouzan exercise)]
\textbf{Q.} Using the extended Euclidean algorithm, find the multiplicative inverses of 7, 132 and 24 in $\mathbb{Z}_{180}$, where they exist.

\textbf{A.} For $b = 7$:
\begin{center}\small
\begin{tabular}{@{}c|c|rrr|rrr@{}}
\toprule
Iteration & $q$ & $r_1$ & $r_2$ & $r$ & $t_1$ & $t_2$ & $t$\\
\midrule
1 & 25 & 180 & 7 & 5 & 0 & 1 & $-25$\\
2 & 1 & 7 & 5 & 2 & 1 & $-25$ & 26\\
3 & 2 & 5 & 2 & 1 & $-25$ & 26 & $-77$\\
4 & 2 & 2 & 1 & 0 & 26 & $-77$ & 180\\
(end) & & \textbf{1} & 0 & & $\mathbf{-77}$ & 180 & \\
\bottomrule
\end{tabular}
\end{center}
$r_1 = 1$, so the inverse exists: $7^{-1} = -77 + 180 = \mathbf{103}$. Check: $7 \times 103 = 721 = 4 \times 180 + 1$ \checkmark.
For $b = 132$ the $r$ columns run $180 = 1 \times 132 + 48$, $132 = 2 \times 48 + 36$, $48 = 1 \times 36 + 12$, $36 = 3 \times 12 + 0$, ending with $r_1 = 12 \ne 1$: \textbf{no inverse}. For $b = 24$: $180 = 7 \times 24 + 12$, $24 = 2 \times 12 + 0$, again $r_1 = 12$: \textbf{no inverse}. Both share the factor 12 with 180.
\end{worked}

\subsection{Linear Diophantine equations \texorpdfstring{$ax + by = c$}{ax + by = c} \hot}
\likely{``Find the particular and general solutions of $ax + by = c$'' (3 marks), as in 2025--26 Q5; occasionally a word problem with non-negative solutions.}
\imp{A linear Diophantine equation in two variables is $ax + by = c$, where $a$, $b$, $c$ are integers and only integer solutions $x$, $y$ are wanted.} Over the reals such an equation has a whole line of solutions; restricting to integers leaves either none or infinitely many, and the extended Euclidean algorithm decides which and finds them all. The three facts below are derived from Bézout's identity, with $d = \gcd(a, b)$.

\textbf{No solution unless $d \mid c$.} For every pair of integers $x, y$, $d$ divides $ax + by$ (Property 4 of \S\ref{sec:gcd}). If $d \nmid c$, no choice of $x$, $y$ can make $ax + by$ equal to $c$. For example $15x + 12y = 13$ has no solution, because $\gcd(15, 12) = 3$ and $3 \nmid 13$.

\textbf{A particular solution when $d \mid c$.} Divide the equation by $d$: $a'x + b'y = c'$ with $a' = a/d$, $b' = b/d$, $c' = c/d$ and $\gcd(a', b') = 1$. The extended algorithm on $(a', b')$ gives $s$ and $t$ with $a's + b't = 1$, and multiplying by $c'$ gives $a'(c's) + b'(c't) = c'$. So $x_0 = (c/d)\,s$, $y_0 = (c/d)\,t$ is a solution. (Running the algorithm on $(a, b)$ directly gives the same $s$, $t$, since $sa + tb = d$ is the same equation multiplied by $d$.)

\textbf{The general solution.} Let $(x, y)$ be any other solution. Subtracting $a'x_0 + b'y_0 = c'$ from $a'x + b'y = c'$ gives $a'(x - x_0) = b'(y_0 - y)$. So $b'$ divides $a'(x - x_0)$, and since $b'$ is co-prime to $a'$, it divides $x - x_0$. (This step, \emph{Euclid's lemma}, follows from Bézout as well: multiplying $a's + b't = 1$ by $m = x - x_0$ gives $m = s\,a'm + t\,b'm$, and $b'$ divides both terms.) Hence $x = x_0 + k b'$ for some integer $k$, and substituting back gives $y = y_0 - k a'$. Conversely every $k$ gives a solution, because $a'(x_0 + kb') + b'(y_0 - ka') = a'x_0 + b'y_0$. These are the slide's formulas:
\[
\text{particular: } x_0 = \frac{c}{d}\,s, \quad y_0 = \frac{c}{d}\,t, \qquad\qquad
\text{general: } x = x_0 + k\,\frac{b}{d}, \quad y = y_0 - k\,\frac{a}{d}, \quad k \in \mathbb{Z}.
\]
The solutions are evenly spaced: $x$ steps by $b/d$ while $y$ steps by $a/d$ in the opposite direction, which keeps $ax + by$ constant. The linear congruence $ax \equiv b \pmod n$ of \S\ref{sec:modular} is the same problem in disguise, since it means $ax + ny = b$ for some integer $y$; the condition $d \mid b$ and the spacing $n/d$ of its $d$ solutions are this result read modulo $n$.

\begin{mustknow}
$s \cdot a + t \cdot b = \gcd(a, b)$: update $r$, $s$, $t$ by the same rule (new $=$ first $- q \times$ second), stop when $r_2 = 0$, read $r_1$, $s_1$, $t_1$.
Inverse: run with $r_1 = n$, $r_2 = b$; if the final $r_1 = 1$, then $b^{-1} = t_1 \bmod n$, otherwise no inverse.
$ax + by = c$: solvable iff $d = \gcd(a, b) \mid c$; $x = \frac{c}{d}s + k\frac{b}{d}$, \;$y = \frac{c}{d}t - k\frac{a}{d}$.
\end{mustknow}

The slide examples. For $21x + 14y = 35$: $d = \gcd(21, 14) = 7$ and $7 \mid 35$, so there are infinitely many solutions. Dividing by 7 gives $3x + 2y = 5$, and the algorithm on $(3, 2)$ gives $s = 1$, $t = -1$. The particular solution is $x_0 = 5 \times 1 = 5$, $y_0 = 5 \times (-1) = -5$, and the general solution is $x = 5 + 2k$, $y = -5 - 3k$: $(5, -5), (7, -8), (9, -11), \dots$, each of which satisfies $21x + 14y = 35$.
For cashing a Rs.~100 cheque in Rs.~20 and Rs.~5 notes, $20x + 5y = 100$: $d = 5$ divides 100, the reduced equation is $4x + y = 20$, and $4s + t = 1$ gives $s = 0$, $t = 1$. The particular solution is $x_0 = 0 \times 20 = 0$, $y_0 = 1 \times 20 = 20$, and the general one is $x = k$, $y = 20 - 4k$. Only non-negative note counts make sense, so $0 \le k \le 5$: $(0, 20)$, $(1, 16)$, $(2, 12)$, $(3, 8)$, $(4, 4)$, $(5, 0)$.

\begin{worked}[: 2025--26 Q5 (3 marks), solve $19x + 13y = 20$]
\textbf{Q.} Find the particular and general solutions of $19x + 13y = 20$.

\textbf{A.} Run the extended Euclidean algorithm on $a = 19$, $b = 13$:
\begin{center}\small
\begin{tabular}{@{}r|rrr|rrr|rrr@{}}
\toprule
$q$ & $r_1$ & $r_2$ & $r$ & $s_1$ & $s_2$ & $s$ & $t_1$ & $t_2$ & $t$\\
\midrule
1 & 19 & 13 & 6 & 1 & 0 & 1 & 0 & 1 & $-1$\\
2 & 13 & 6 & 1 & 0 & 1 & $-2$ & 1 & $-1$ & 3\\
6 & 6 & 1 & 0 & 1 & $-2$ & 13 & $-1$ & 3 & $-19$\\
 & \textbf{1} & 0 & & $\mathbf{-2}$ & 13 & & \textbf{3} & $-19$ & \\
\bottomrule
\end{tabular}
\end{center}
$d = \gcd(19, 13) = 1$, $s = -2$, $t = 3$ (check: $-2 \times 19 + 3 \times 13 = -38 + 39 = 1$). Since $1 \mid 20$, solutions exist.

\textbf{Particular:} $x_0 = (c/d)\,s = 20 \times (-2) = \mathbf{-40}$, \;$y_0 = (c/d)\,t = 20 \times 3 = \mathbf{60}$. \;Check: $19(-40) + 13(60) = -760 + 780 = 20$ \checkmark.

\textbf{General:} $x = -40 + 13k$, \;$y = 60 - 19k$, \;$k \in \mathbb{Z}$.
\begin{center}\small
\begin{tabular}{@{}c|cccccc@{}}
\toprule
$k$ & 0 & 1 & 2 & 3 & 4 & 5\\
\midrule
$(x, y)$ & $(-40, 60)$ & $(-27, 41)$ & $(-14, 22)$ & $\mathbf{(-1, 3)}$ & $(12, -16)$ & $(25, -35)$\\
\bottomrule
\end{tabular}
\end{center}
The smallest solution is $(-1, 3)$ ($k = 3$), so the general solution can also be written $x = -1 + 13k$, $y = 3 - 19k$. Check: $19(-1) + 13(3) = -19 + 39 = 20$ \checkmark.
\end{worked}

\begin{worked}[: non-negative solutions (Forouzan-style word problem)]
\textbf{Q.} A post office sells only 39-cent and 15-cent stamps. Find the numbers of stamps a customer can buy to pay exactly \$2.70 of postage.

\textbf{A.} Let $x$ and $y$ be the numbers of 39-cent and 15-cent stamps: $39x + 15y = 270$, with $x, y \ge 0$. $d = \gcd(39, 15) = 3$ ($39 = 2 \times 15 + 9$, $15 = 1 \times 9 + 6$, $9 = 1 \times 6 + 3$, $6 = 2 \times 3$), and $3 \mid 270$. Dividing by 3: $13x + 5y = 90$.
\begin{center}\small
\begin{tabular}{@{}r|rrr|rrr|rrr@{}}
\toprule
$q$ & $r_1$ & $r_2$ & $r$ & $s_1$ & $s_2$ & $s$ & $t_1$ & $t_2$ & $t$\\
\midrule
2 & 13 & 5 & 3 & 1 & 0 & 1 & 0 & 1 & $-2$\\
1 & 5 & 3 & 2 & 0 & 1 & $-1$ & 1 & $-2$ & 3\\
1 & 3 & 2 & 1 & 1 & $-1$ & 2 & $-2$ & 3 & $-5$\\
2 & 2 & 1 & 0 & $-1$ & 2 & $-5$ & 3 & $-5$ & 13\\
 & \textbf{1} & 0 & & \textbf{2} & $-5$ & & $\mathbf{-5}$ & 13 & \\
\bottomrule
\end{tabular}
\end{center}
$s = 2$, $t = -5$ ($2 \times 13 - 5 \times 5 = 1$). Particular: $x_0 = 90 \times 2 = 180$, $y_0 = 90 \times (-5) = -450$. General: $x = 180 + 5k$, $y = -450 - 13k$.
Non-negative counts need $180 + 5k \ge 0$, i.e.\ $k \ge -36$, and $-450 - 13k \ge 0$, i.e.\ $k \le -34.6$, so $k \in \{-36, -35\}$.
\textbf{Answer: $(x, y) = (0, 18)$ or $(5, 5)$}: eighteen 15-cent stamps, or five of each. Check: $15 \times 18 = 270$ and $39 \times 5 + 15 \times 5 = 195 + 75 = 270$ \checkmark.
\end{worked}

\begin{trap}[: where marks are lost]
\begin{itemize}[leftmargin=1.5em]
  \item For an inverse, put the \textbf{modulus first} ($r_1 = n$). Starting with $r_1 = b$ makes the inverse appear in the $s$ column, which the $t$-only table has dropped.
  \item Read the answer from the final row's \emph{first} columns ($r_1$, $s_1$, $t_1$), not from $t_2$ (which is $\pm n$) or from the last computed $t$.
  \item A negative $t_1$ is not an error: add $n$ (in RSA, add $\phi(n)$, never $n$). Check $b \times b^{-1} \bmod n = 1$ before writing the answer.
  \item If the final $r_1 \ne 1$, write ``no inverse, since $\gcd = r_1$''; do not report $t_1$.
  \item Show the full table with headed columns; the marks are awarded per row, and the ``iteration $k$'' sub-questions can only be answered from it.
\end{itemize}
\end{trap}

\begin{fixbox}
\begin{itemize}
  \item All slide tables are correct (recomputed by script): $(161, 28) \to 7, -1, 6$; $(17, 0) \to 17, 1, 0$; $(0, 45) \to 45, 0, 1$; $(84, 320) \to 4, -19, 5$; $11^{-1} \bmod 26 = 19$; $23^{-1} \bmod 100 = 87$; $12^{-1} \bmod 26$ does not exist; $21x + 14y = 35$ and the cheque solutions.
  \item Lecture 4, Diophantine slide: ``if $d \nmid c$, the equation has no solution'' has the ``not'' stroke hand-drawn on the slide. A text copy of the older slides lost it and reads ``if $d \mid c$, the equation has no solution'', which is wrong; the condition for \emph{no} solution is $d \nmid c$.
\end{itemize}
\end{fixbox}
